\uuid{E3SD}
\titre{Raisonnement}
\auteur{bodin ; hanani ; mzari}
\organisation{exo7}
\licence{CC BY-NC-SA 4.0}
\datecreate{2020-01-20}
\niveau{L1}
\module{Algèbre}
\chapitre{Logique, ensemble, raisonnement}
\sousChapitre{Absurde et contraposée}
\difficulte{3}
\theme{Raisonnement}
\exo7qcmid{lille-1.logique.030}

\contenu{
\begin{question}
\qclassification{100.03,100.04}
\qtags{motcle=raisonnement}
Je veux montrer que $\sqrt{13} \notin \Qq$ par un raisonnement par l'absurde. Quel schéma de raisonnement est adapté ?
\begin{answers}
    \good{Je suppose que $\sqrt{13}$ est rationnel et je cherche une contradiction.}

    \bad{Je suppose que $\sqrt{13}$ est irrationnel et je cherche une contradiction.}

    \bad{J'écris $13 = \frac{p}{q}$ (avec $p,q$ entiers) et je cherche une contradiction.}

    \good{J'écris $\sqrt{13} = \frac{p}{q}$ (avec $p,q$ entiers) et je cherche une contradiction.}
\end{answers}
\begin{explanations}
Par l'absurde on suppose que $\sqrt{13} \in \Qq$, c'est-à-dire que c'est un nombre rationnel, autrement dit qu'il s'écrit $\frac{p}{q}$, avec $p$, $q$ entiers. Voir la preuve que $\sqrt{2} \notin \Qq$.
\end{explanations}
\end{question}
}
